\section{Comparison of the supplied constructions}

The comparison was performed against ProveIt revision
\file{dc523e7b06ffb414093c90aa035500ca283e5f75}, the requested PA209
revision \file{69d98241a721ead20b4f0b17df9062950b40fbfe}, and the
published PA196 table~\cite{pa196,pa209,incoming}. The latter's SHA-256 is
\file{e8bbccd4cfa325251ca95066db116b5aa766e404b0971c33c02772be1c3617a8}.
State counts throughout count working controls; the distinguished halt
control is separate. Each machine is binary, has one two-way infinite tape,
starts from all zeroes, and moves one cell left or right on every transition.

The classification and the earlier reductions described in this section
already occur in the incoming research. They have been freshly replayed
and their input bytes rebound to the current original ZIP members; they
are not new state reductions discovered in this investigation.

\subsection{The full U208 collection has twenty graph classes}

The nine headline tables in the five original U208 packages form exactly
two rooted transition-graph isomorphism classes, called A and B. Each
isomorphism preserves the initial and halting controls, read symbols,
written symbols, and movement directions. Raw byte hashes differ under
renumbering, so comparing hashes alone would overcount independent designs.

\begin{center}\small
\begin{tabular}{lll}
\hline
Package, dated 29 October 2026 & Principal table & Alternative\\
\hline
\texttt{U208\_ConPA} & A & ---\\
\texttt{U208\_ConPA\_Conditional\_Context} & A & B\\
\texttt{U208\_ConPA\_Entry\_Context} & A & B\\
\texttt{U208\_ConPA\_History} & A & B\\
\texttt{u208\_research} & B & A\\
\hline
\end{tabular}
\end{center}

The complete collection is larger. The archive
\texttt{pa208\_research\_2026-10-29.zip} contains eighteen additional
accepted 208-state helper constructions, nine based on each class.
All eighteen graphs are pairwise nonisomorphic and differ from A and B.
Consequently the accepted U208 material audited here consists of
\emph{27 presentations in 20 rooted graph classes}. Each new helper
construction deletes one old state but introduces one fresh helper, so
its net state count stays 208.

The retained comparison checks 30 selected literal inputs: these 27 U208
tables, Chang's PA209 and PA196, and the previously supplied alternative196.
It reconstructs all sixteen isomorphism maps among the nine headline
tables and verifies every transition under each map. It also independently
rechecks the exact patches, incoming-edge predicates, and positive local
whole-tape simulations for all eighteen helper variants. All 30 table
hashes were checked directly against their current original archive
members or repository files, rather than only against copied inputs.
The portable results are in \texttt{comparison/verified\_comparison.json}.

\subsection{Source sizes and physical sizes}

\begin{center}\small
\begin{tabular}{lrrrrrr}
\hline
Design & States & Sites & Internal nodes & Leaf types & DAG total & Max. path\\
\hline
U208 A &208&3909&174&23&197&42\\
U208 B &208&3909&174&23&197&41\\
Chang PA209 &209&3948&173&23&196&33\\
Chang PA196 &196&3812&159&23&182&38\\
Incoming alternative196 &196&3812&159&23&182&36\\
\hline
\end{tabular}
\end{center}

A source site is a primitive instruction occurrence. Internal nodes count
distinct binary subtree pairs, while the total DAG also counts distinct
primitive instructions and their operands. Thus U208's total DAG count
197 must not be compared directly with PA196's internal-node count 159.

A and B have the same ordered primitive word and the same exact ordinal
control-flow graph; 720 binary instruction addresses differ. Chang's
PA209 is distinct from the U209 parents inside the U208 reports. Sharing
an earlier tape engine does not identify their literal state labels or
make their source optimizations additive.

The PA196 budget is
\[
36\text{ engine states}+13\text{ leaf routers}
 +147\text{ source readers}=196.
\]
The 147 readers implement 159 internal source nodes. In particular, the
six contextual identifications in the axiom body are already included
in 196; they implement 68 axiom-grammar nodes using 62 physical readers.

The replay also checks the complete finite initializations, not merely
a tape prefix. PA209 reaches its clean frame after 58,210 transitions;
PA196 does so after 58,154. Their frames agree up to the documented
control-label correspondence: the head is at position $-2$, the register
file represents candidate 2, the gap is 48, and there are 26 spare blocks
of value 1. These are initialization facts, not evidence of nonhalting.

\subsection{Why the final U208 deletion does not directly transfer}

In the Joint native U209 numbering, the final U208 reduction identifies
roles 172 and 211. An entry to 172 has predecessor bits $10$; when 211
scans zero, its immediately preceding bit is one. Re-reading that history
lets the merged control select the correct continuation. The replacement
blocks for old state/read pairs
$(172,0),(172,1),(211,0),(211,1)$ have lengths $7,1,3,1$.
Their endpoints agree on the entire tape and head position, not just on a
finite numerical register sample. Thus the exact simulation clock is
\[
 T_{208}=T_{209}+6N_{172,0}+2N_{211,0},
\]
where the $N$'s count visits in the parent execution. The incoming-edge
predicates and positive block lengths support a PA-formalizable finite
stuttering argument, independently of the underlying theory calibration.

PA209-to-PA196 changes the axiom emitter and its coding. It transposes the
arguments of the universal-quantifier Cantor pair, removes an old axiom
head, adds guarded equality vacuity, and changes formula and emitter
ordering. In the native comparison, 75 old controls disappear and 62 new
ones are introduced: $209-75+62=196$. Roles 172 and 211 are already absent.
Moreover, a complete finite scan of PA196 finds no state $p$ whose two
rows have the form
\[
(p,0)\mapsto(0,L,q),\qquad(p,1)\mapsto(1,R,q)
\]
for the same working state $q$. Therefore no state renaming transplants
that exact U208 mixed-entry motif. This does not exclude more general
contextual sharing or a different source allocation.

\subsection{The arithmetic endpoint must be stated accurately}

The upstream theorem inventory exports the internal statement
\[
\mathsf{PA}\vdash\neg\operatorname{Halt}(M_{196})
   \longrightarrow\operatorname{Con}(\mathsf{PA}),
\]
and external Lean-metatheoretic nonhalting and independence results.
It does not export the full internal biconditional. This is a statement
about the supplied declarations, not a proof that the missing converse
is false or unprovable. A uniform syntactic extraction of a PA refutation
from a halted literal trace can establish the converse; external
HF-model soundness alone does not supply that PA-internal theorem.
The incoming packets contain additional written proof-translation
arguments with explicitly stated implementation boundaries. None of the
state-count experiments below fills that boundary by itself.

\section{New combinations of previously separate reductions}

The first new experiment composes previously supplied source changes,
then reconstructs and checks their physical implementations. The changes
are: the distributed-separation emitter; the diagonal-vacuity emitter;
the shorter negation block; and the MP-success branch that returns its
consequent in register $H$. These are combined as actual source trees,
not by subtracting their separate state-count estimates.

The negation replacement changes a 47-instruction block
\texttt{push; incT; cons} to a 20-instruction block
\texttt{drainTtoX; incT; pair4}. Its mathematical equivalence requires
the inherited scratch-zero entry invariant. The MP change removes 23
primitive instructions. These logical/source obligations are separate
from the checks that a tape implements the resulting source program.

\begin{center}\small
\begin{tabular}{lrrrr}
\hline
Emitter and changes & Sites & Raw states & Final states & Symbolic cases\\
\hline
Original &3812&202&196&4959\\
Original + MP &3789&202&196&4929\\
Original + short negation &3596&202&197&4671\\
Original + short negation + MP &3573&202&197&4641\\
Distributed &3773&201&198&4908\\
Distributed + MP &3750&201&198&4878\\
Distributed + short negation &3611&202&199&4692\\
Distributed + short negation + MP &3588&202&199&4662\\
Diagonal &3700&203&199&4815\\
Diagonal + MP &3677&203&199&4785\\
Diagonal + short negation &3484&203&200&4527\\
Diagonal + short negation + MP &3461&203&200&4497\\
\hline
\end{tabular}
\end{center}

Here ``raw'' refers to the straightforward allocation before contextual
reader folds. For each source, 200 deterministic seeded fold orders were
examined. The selected complete table passes every symbolic primitive
contract, preserves the 42 explicitly checked engine rows needed by the
register-block interface, and reaches the exact clean blank-start frame
after 58,154 transitions. These are executable symbolic checks, not
finite sampling of register magnitudes. They are also not a new
proof-assistant kernel certificate or an exported PA refutation compiler.

The original and original-plus-MP rows retain 196 states; none is a
strict improvement. The table gives a concrete reason the separate
savings cannot simply be added. Shortening negation removes source work
but disrupts useful contextual sharing, costing a physical state in
these allocations. The reproducible subset is in
\texttt{combinations/}; its replay writes fresh evidence into a separate
\texttt{replayed/} directory.

\subsection{Combining distribution with diagonal vacuity}

A further source experiment combines the two logical emitter changes
with one another. Write $I$ for implication, $U_b$ for universal
quantification, and $N(X)=I(X,1)$. For parameter codes $a,b,c$, set
\[
\begin{split}
 D&=I(c,U_c(c)),\\
 V&=I(M(a,b),U_c(M(a,b))),\\
 P&=I(c,M(c,b)),\\
 L&=I(M(a,b),P),\\
 R&=I(M(c,b),c).
\end{split}
\]
The candidate gate is
\[
 N\!\left(I\!\left(D,
 I\!\left(V,U_c\bigl(I(U_b(L),I(U_b(R),N(P)))\bigr)\right)\right)\right).
\]
It combines the earlier diagonal clause with distribution of the
quantifier over the separation conjunction, using the same original
freshness guard. The preceding head order and one antecedent exchange
are the previously checked distributed-emitter choices.

Exact macro expansion gives 189 tokens, compared with 201 in PA196,
and 3661 source sites. Adding the MP rewrite gives 3638. Applying all
six applicable shorter-negation replacements gives 3499 sites, or 3476
with MP as well. These are exact source-size calculations; they do not
establish a 189-state or a sub-196-state machine. A physical layout must
retain the decoder-entry predicates, maximum-zero-run constraints,
startup, and full symbolic runtime contracts. The source combination
also inherits the outstanding uniform proof-translation and internal
endpoint obligations of its constituent logical changes.

The joint word admits a concrete variable-arity dictionary whose ordinary
binary concatenation cost is 66 after accounting for an already available
two-increment block. A 60-second mixed-integer search supplied this
incumbent; it did not prove optimality. The full witness is retained in
\texttt{combinations/joint\_dictionary/grammar66.json}.
It is not an accepted physical layout. In that fixed dictionary, the
first-part chain through production 208 forces a run of at least seven
consecutive zero path bits: three come from the existing copy-$A$ block,
and four from mandatory left-child descents. The inherited reader bound
is six. The independent \texttt{check\_obstruction.py} verifies the exact
189-token word, cost 66, use of the offending production, and the
seven-zero lower bound without rerunning a solver or trusting a random
association. Thus this particular dictionary cannot be repaired merely
by selecting a different association of its fixed ordered productions.
Changing the productions or the reader remains possible. This local
obstruction is not a lower bound for all 189-token emitter layouts, much
less for all 195-state consistency machines.
